\section{Hilos y Procesos}

Un proceso engloba dos características independientes:
\begin{enumerate}
  \item Propiedad de recursos: Un proceso cuenta con un espacio de direcciones virtual que alberga su imagen (el programa, los datos, la pila y el bloque de control del proceso). Además, tiene asignados recursos como memoria, dispositivos de entrada/salida, canales y archivos. El sistema operativo se encarga de proteger estos recursos para evitar interferencias por parte de otros procesos.
  \item Planificación y ejecución: Un proceso sigue una ruta de ejecución a través de uno o más programas. Posee un estado de ejecución (en ejecución, listo, etc.) y un nivel de prioridad, lo que permite al sistema operativo planificarlo y despacharlo.
\end{enumerate}

Dado que estas características son independientes, muchos sistemas operativos las manejan por separado. A la unidad de ejecución se la conoce como hilo (o proceso ligero), mientras que el proceso continúa siendo la unidad de propiedad de los recursos.

\subsection{Multihilo}

El multihilo es la capacidad de un sistema operativo para gestionar varios hilos que se ejecutan de manera concurrente dentro de un mismo proceso. En este modelo, el proceso actúa como la unidad de asignación y protección de recursos, mientras que los hilos son las unidades de ejecución.

Un proceso aporta:
\begin{itemize}
  \item Un espacio de direcciones virtual que contiene la imagen del proceso.
  \item Acceso protegido a procesadores, archivos, recursos de entrada/salida y otros procesos.
\end{itemize}

Cada hilo dispone de:
\begin{itemize}
  \item Su propio estado de ejecución y contexto guardado.
  \item Su propia pila de ejecución y almacenamiento local.
  \item Acceso a la memoria y los recursos compartidos de su proceso contenedor.
\end{itemize}

En un proceso multihilo, todos los hilos comparten el mismo espacio de direcciones y los recursos generales, pero cada uno conserva de forma individual su bloque de control, sus registros, su pila y su estado de ejecución.

Las principales ventajas del uso de hilos son:
\begin{enumerate}
  \item La creación de un hilo es considerablemente más rápida que la de un proceso.
  \item La finalización de un hilo es más veloz que la de un proceso.
  \item El cambio de contexto entre hilos requiere menos recursos informáticos que entre procesos.
  \item La comunicación entre hilos es más eficiente, ya que comparten la misma memoria y recursos sin necesidad de intervenir el sistema operativo.
\end{enumerate}

Por lo general, esto hace que las aplicaciones conformadas por unidades de ejecución relacionadas funcionen de manera mucho más eficiente si se implementan con hilos en lugar de procesos independientes. Esto es un importante detalle, ya que para que una aplicación aproveche el multihilo, debe cumplir con dos condiciones: \texttt{(i)} tener una lógica paralelizable y \texttt{(ii)} estar programada de forma explícita para invocar las funciones de concurrencia del sistema operativo.

En sistemas multiprocesador, varios hilos de un mismo proceso pueden ejecutarse de forma simultánea, mientras que en sistemas con un solo procesador el uso de hilos facilita la lógica de los programas que realizan distintas tareas independientes.

Los usos habituales de los hilos incluyen:
\begin{itemize}
  \item Separar las tareas en primer plano de las de segundo plano.
  \item Realizar operaciones asíncronas, como copias de seguridad periódicas.
  \item Superponer el procesamiento de datos con las operaciones de entrada/salida para mejorar el rendimiento global.
  \item Simplificar el diseño de aplicaciones modulares.
\end{itemize}

La figura \ref{fig:modelo_hilos} representa la comparativa entre el modelo de proceso anteriormente visto (\emph{Single-Threaded process model}) y el modelo utilizando hilos y procesos (\emph{Multithreaded process model}).

\begin{figure}[!ht]
  \centering
  \includegraphics[width=0.9\textwidth]{chapters/02_unit/media/threads.pdf}
  \caption{Modelos de un solo hilo y múltiples hilos.}
  \label{fig:modelo_hilos}
\end{figure}

En los sistemas operativos que admiten multihilo, la planificación y el despacho se realizan a nivel de hilo, por lo que toda la información vinculada a la ejecución se guarda en estructuras de control propias del hilo. No obstante, las operaciones que afectan al proceso en su totalidad ---como suspenderlo o finalizarlo--- repercuten en todos sus hilos, puesto que comparten el mismo espacio de direcciones y recursos.

\subsection{Funcionamiento de los hilos}

Al igual que los procesos, los hilos poseen estados de ejecución y pueden sincronizarse entre sí.

\subsubsection{Estados de los hilos}

Los estados principales de un hilo son \textbf{Ejecutando} (Running), \textbf{Listo} (Ready) y \textbf{Bloqueado} (Blocked). En general, los estados de suspensión (\textit{Suspend}) se consideran conceptos propios del proceso, ya que si un proceso es suspendido o intercambiado fuera de memoria, todos sus hilos también lo son al compartir el mismo espacio de direcciones.

Las principales operaciones sobre un hilo son:
\begin{enumerate}
  \item \textbf{Creación (Spawn):} se crea un nuevo hilo con su propio contexto de registros y pila, quedando en estado \textbf{Listo}.
  \item \textbf{Bloqueo (Block):} el hilo espera un evento, guardando su contexto para reanudar la ejecución posteriormente.
  \item \textbf{Desbloqueo (Unblock):} cuando ocurre el evento esperado, el hilo vuelve al estado \textbf{Listo}.
  \item \textbf{Finalización (Finish):} al terminar su ejecución, se liberan su contexto y su pila.
\end{enumerate}

Un aspecto importante es determinar si el bloqueo de un hilo implica el bloqueo de todo el proceso. Si esto ocurre, los demás hilos del proceso no pueden ejecutarse, reduciendo las ventajas del multihilo. En los sistemas donde los hilos pueden bloquearse de forma independiente, mientras un hilo espera un evento, otros hilos del mismo proceso continúan ejecutándose o permanecen listos para hacerlo.

En un sistema con un único procesador, varios hilos pertenecientes a distintos procesos pueden ejecutarse de forma intercalada. El cambio de ejecución ocurre cuando un hilo se bloquea o cuando finaliza su límite de tiempo (\textit{quantum}). En un sistema operativo moderno (como Linux o Windows), la cola de planificación está llena de hilos, no de procesos. Cuando el procesador termina su \textit{quantum} o un hilo se bloquea, el planificador busca el siguiente hilo en estado \textbf{Listo (Ready)} para ejecutarlo en la CPU.

Al realizar el cambio, el kernel compara a qué proceso pertenecen el hilo saliente y el entrante, pudiendo darse dos situaciones:
\begin{itemize}
  \item \textbf{Cambio intra-proceso (ligero):} Si ambos hilos pertenecen al mismo proceso (ej. Hilo 1 del Proceso A a Hilo 2 del Proceso A), el kernel solo conmuta el contexto del procesador: contador de programa, registros de propósito general y puntero de pila. El registro que controla la memoria virtual se mantiene igual, por lo que es una operación extremadamente rápida.

  \item \textbf{Cambio inter-proceso (pesado):} Si los hilos pertenecen a procesos distintos (ej. Hilo 2 del Proceso A a Hilo 1 del Proceso B), además de conmutar registros y pila, el sistema operativo debe realizar un cambio completo de contexto de memoria para cargar las tablas de páginas del nuevo proceso y debe invalidar o re-etiquetar la caché de traducción de direcciones (\textit{Translation Lookaside Buffer}, \textbf{TLB}), lo cual implica un costo significativo de rendimiento.
\end{itemize}

Por lo tanto, si bien la CPU reparte su tiempo de forma intercalada entre hilos de distintos procesos, cada vez que esa intercalación cruza la frontera de un proceso, el kernel debe efectuar un cambio de contexto completo, perdiendo la ventaja de velocidad del cambio ligero intra-proceso.

\subsubsection{Sincronización de hilos}

Todos los hilos de un proceso comparten el mismo espacio de direcciones y recursos, como archivos abiertos. Por ello, las modificaciones realizadas por un hilo son visibles para los demás, lo que hace necesaria la sincronización para evitar interferencias y corrupción de datos.

Los problemas y mecanismos utilizados para sincronizar hilos son, en esencia, los mismos que se emplean para la sincronización entre procesos y se estudiarán con mayor detalle en capítulos posteriores.

\section{Tipos de hilos}

\subsection{Hilos de usuario y hilos de kernel}

Existen dos formas principales de implementar hilos: los \textbf{hilos de usuario} (\textit{User-Level Threads, ULT}) y los \textbf{hilos de kernel} (\textit{Kernel-Level Threads, KLT}), también conocidos como \textit{kernel-supported threads} o \textit{lightweight processes}.

\begin{figure}[ht]
  \centering
  \includegraphics[width=0.9\textwidth]{chapters/02_unit/media/types_of_threads.pdf}
  \caption{Hilos de usuario e Hilos de Kernel.}
  \label{fig:usr_and_kernel_threads}
\end{figure}

\subsubsection{Hilos de usuario}

En una implementación pura de ULT, toda la gestión de los hilos se realiza en espacio de usuario mediante una \textbf{biblioteca de hilos} (\textit{threads library}), mientras que el kernel desconoce la existencia de dichos hilos y únicamente administra el proceso.

La biblioteca de hilos proporciona las funciones necesarias para:
\begin{itemize}
  \item Crear y destruir hilos.
  \item Planificar su ejecución.
  \item Intercambiar información entre hilos.
  \item Guardar y restaurar el contexto de ejecución.
\end{itemize}
Cuando un proceso crea un nuevo hilo, la biblioteca le asigna su propio contexto (registros, contador de programa y pila) y lo coloca en estado \textbf{Ready}. Los cambios de contexto entre hilos también son realizados por la biblioteca, sin intervención del kernel.

Como el kernel solo conoce al proceso, este mantiene un único estado de ejecución (\textbf{Running}, \textbf{Ready}, \textbf{Blocked}, etc.). Por ello, el planificador del sistema operativo decide qué proceso ejecutar, mientras que la biblioteca decide cuál de sus hilos utilizar.

Esta organización produce algunas consecuencias importantes:
\begin{itemize}
  \item Si un hilo realiza una llamada al sistema bloqueante, el kernel bloquea al proceso completo, impidiendo la ejecución de todos sus hilos.
  \item Si el proceso es interrumpido o pierde el procesador, todos sus hilos dejan de ejecutarse hasta que el proceso vuelva a ser planificado.
  \item Un hilo puede bloquearse mientras otro hilo del mismo proceso continúa ejecutándose, siempre que el bloqueo sea administrado por la biblioteca de hilos y no por el kernel.
\end{itemize}
Las principales ventajas de los ULT son:
\begin{enumerate}
  \item Los cambios de contexto entre hilos son rápidos, ya que no requieren cambiar al modo kernel.
  \item Cada aplicación puede utilizar su propio algoritmo de planificación de hilos.
  \item Pueden implementarse sobre cualquier sistema operativo sin modificar el kernel.
\end{enumerate}
Sus principales desventajas son:
\begin{enumerate}
  \item Una llamada al sistema bloqueante detiene a todos los hilos del proceso.
  \item El proceso solo puede ejecutarse en un procesador a la vez, por lo que los hilos no pueden aprovechar el paralelismo de sistemas multiprocesador.
\end{enumerate}

Existen técnicas para mitigar estas limitaciones. Una de ellas es el \textbf{jacketing}, que consiste en reemplazar una llamada bloqueante por una rutina de la biblioteca que verifica si la operación puede realizarse inmediatamente. Si el recurso no está disponible, el hilo se bloquea a nivel de usuario y otro hilo del mismo proceso continúa ejecutándose.

\subsubsection{Hilos de kernel}

En una implementación pura de \textbf{Kernel-Level Threads (KLT)}, toda la gestión de los hilos es realizada por el kernel. Las aplicaciones únicamente utilizan una interfaz (\textit{API}) para acceder a los servicios de administración de hilos.

El kernel mantiene información tanto del proceso como de cada uno de sus hilos y realiza la planificación a nivel de hilo. Esto elimina las principales limitaciones de los ULT:
\begin{itemize}
  \item Varios hilos del mismo proceso pueden ejecutarse simultáneamente en diferentes procesadores.
  \item Si un hilo se bloquea, el kernel puede planificar otro hilo del mismo proceso.
  \item El propio kernel puede implementar sus rutinas utilizando múltiples hilos.
\end{itemize}
La principal desventaja es que cada cambio de contexto entre hilos requiere la intervención del kernel, lo que implica un cambio de modo (\textit{user mode} $\leftrightarrow$ \textit{kernel mode}) y un mayor costo que en los ULT. No obstante, el rendimiento obtenido depende de la aplicación: si los hilos realizan frecuentes llamadas al sistema, la diferencia de rendimiento entre ULT y KLT puede ser reducida.

\subsubsection{Enfoques combinados}

Como ya pudo observar en la figura \ref{fig:usr_and_kernel_threads}, hay una alternativa adicional. Algunos sistemas operativos utilizan un enfoque híbrido que combina ULT y KLT. En este modelo, la creación, planificación y sincronización de los hilos se realiza principalmente en espacio de usuario, mientras que los hilos de usuario se asignan a uno o más hilos de kernel.

Este enfoque permite que:
\begin{itemize}
  \item Múltiples hilos de una misma aplicación se ejecuten en paralelo sobre distintos procesadores.
  \item Una llamada al sistema bloqueante no detenga a todos los hilos del proceso.
  \item Se combinen la eficiencia de los ULT con las capacidades de planificación y paralelismo de los KLT.
\end{itemize}
Si está bien diseñado, un sistema híbrido aprovecha las ventajas de ambos modelos y reduce sus desventajas.
